\renewcommand{\arraystretch}{1.35}
\captionsetup[longtable]{justification=raggedright,singlelinecheck=false}

\begin{longtable}{|p{3.8cm}|p{5cm}|p{4cm}|}
\caption{Pemetaan Tahap Metodologi}
\label{tbl:mapping_methodology_requirements} \\
\hline
\textbf{Tahap Metodologi} &
\textbf{Aktivitas Utama} &
\textbf{Kebutuhan yang Dipenuhi (F)} \\
\hline
\endfirsthead

\multicolumn{3}{l}{Tabel \ref{tbl:mapping_methodology_requirements} Pemetaan Tahap Metodologi (lanjutan)} \\
\hline
\textbf{Tahap Metodologi} &
\textbf{Aktivitas Utama} &
\textbf{Kebutuhan yang Dipenuhi (F)} \\
\hline
\endhead

\hline
\multicolumn{3}{r}{\textit{bersambung ke halaman berikutnya}} \\
\endfoot

\hline
\endlastfoot

% =================== DATA PREPARATION ======================
\multicolumn{3}{|c|}{\textbf{Data Preparation}} \\ \hline

\textit{Data Selection \& Cleaning} &
Identifikasi objek target dari spesifikasi OpenAPI Kubernetes, seleksi blok \textit{definitions}, dan pembersihan tipe generik yang tidak relevan untuk pembuatan YAML. &
F-01 (\textit{Structured Object Representation}),
F-05 (\textit{Semantic Embedding}) \\ \hline

Konstruksi Topologi Graf &
Transformasi struktur JSON menjadi node–edge dengan relasi referensial dan hierarkis. &
F-01 (\textit{Structured Object Representation}),
F-02 (\textit{Structured Relation Retrieval}),
F-05 (\textit{Semantic Embedding}) \\ \hline

\textit{Semantic Vectorization} &
Ekstraksi deskripsi objek menjadi \textit{embedding} numerik untuk pencarian semantik. &
F-05 (\textit{Semantic Embedding}) \\ \hline

Neo4j \textit{Data Loading} &
Memuat hasil konstruksi graf ke basis data Neo4j agar dapat ditelusuri kembali saat retrieval. &
F-01 (\textit{Structured Object Representation}) \\ \hline


% =================== MODELLING ======================
\multicolumn{3}{|c|}{\textbf{Modelling}} \\ \hline

\textit{Intent Classification Modeling} &
Klasifikasi pertanyaan pengguna ke dalam tipe \textit{intent} untuk menentukan strategi dan kedalaman penelusuran graf. &
F-03 (\textit{Intent Classification}) \\ \hline

\textit{Retrieval Strategy Design} &
Perancangan traversal graf berdasarkan entitas target, relasi, dan kebutuhan pengguna. &
F-02 (\textit{Structured Relation Retrieval}),
F-04 (\textit{Context Construction \& Injection}) \\ \hline

\textit{Prompt Engineering} &
Perancangan template \textit{prompt} untuk dua peran LLM serta injeksi konteks dan pengaturan format keluaran. &
F-04 (\textit{Context Construction \& Injection}),
F-07 (\textit{Conditional YAML Generation Constraint}) \\ \hline

\textit{Validation Logic} &
Integrasi prosedur validasi YAML bertingkat sebagai mekanisme \textit{fail-safe} yang memastikan keluaran memenuhi standar sintaksis, skema, dan kelengkapan \textit{field}. &
F-07 (\textit{Conditional YAML Generation Constraint}) \\ \hline

\textit{Conversation Memory Management} &
Penyimpanan dan pengambilan riwayat percakapan dalam \textit{pipeline} LangGraph untuk mendukung konteks \textit{multi-turn}. &
F-03 (\textit{Intent Classification}),
F-04 (\textit{Context Construction \& Injection}) \\ \hline


% =================== INPUT ======================
\multicolumn{3}{|c|}{\textbf{Input Processing}} \\ \hline

Pertanyaan Pengguna &
Input teks berupa permintaan teknis Kubernetes dalam bahasa alami yang diajukan pengguna melalui antarmuka percakapan. &
F-03 (\textit{Intent Classification}) \\ \hline


% =================== OUTPUT ======================
\multicolumn{3}{|c|}{\textbf{Output Generation}} \\ \hline

Konten Jawaban &
Keluaran sistem berupa narasi penjelas atau blok YAML konfigurasi Kubernetes sesuai permintaan pengguna. &
F-07 (\textit{Conditional YAML Generation Constraint}) \\ \hline

\textit{Graph Traceability} &
Penelusuran asal-usul \textit{field} melalui relasi \textit{node–edge} di dalam graf. &
F-02 (\textit{Structured Relation Retrieval}) \\ \hline


% =================== EVALUATION ======================
\multicolumn{3}{|c|}{\textbf{Evaluation}} \\ \hline

\textit{Answer Quality} (AnsQ) &
Evaluasi kualitas jawaban (\textit{Answer Quality}); definisi metrik lengkap pada Subbab~\ref{sec:metrik-evaluasi}. &
-- (Metrik domain, bukan requirement fungsional) \\ \hline

\textit{Retrieval Quality} (RetQ) &
Evaluasi kualitas \textit{retrieval} (\textit{Retrieval Quality}); definisi metrik lengkap pada Subbab~\ref{sec:metrik-evaluasi}. &
-- (Metrik domain, bukan requirement fungsional) \\ \hline

\textit{Reasoning Quality} (ReaQ) &
Evaluasi kualitas penalaran (\textit{Reasoning Quality}); definisi metrik lengkap pada Subbab~\ref{sec:metrik-evaluasi}. &
-- (Metrik domain, bukan requirement fungsional) \\ \hline


% =================== DEPLOYMENT ======================
\multicolumn{3}{|c|}{\textbf{Deployment}} \\ \hline

Antarmuka Web Interaktif (Streamlit) &
Penyajian sistem melalui antarmuka web berbasis Streamlit (\texttt{main.py}) yang menampilkan \textit{chat interface}, jawaban LLM, blok YAML (jika relevan), serta visualisasi \textit{Retrieval Trace} berbasis Graphviz. &
F-06 (\textit{Web Interaction Interface}) \\ \hline

\end{longtable}
